
This chapter closes the dissertation. Section 9.1 states how each objective was met and what the two research questions were answered with, \S9.2 sets out the contributions and who can use them, \S9.3 states the limitations honestly and turns each into the work that would resolve it, and \S9.4 is a personal reflection on how the project was actually conducted.

\section{Summary of the dissertation}

The aim was to determine whether spatial-relationship annotation for robot-acquired images can be automated, and whether the resulting labels are good enough to replace the human ones. Five objectives carried that aim, and each is met with evidence rather than assertion.

\textbf{O1, build.} A fully-automatic pipeline annotates all 836 images in about five minutes on a 6 GB consumer GPU, with no human deciding any label, and writes the dataset's own formats byte-compatibly (Chapter 3). Verified by a load-write round trip reproducing boxes and labels with zero error.

\textbf{O2, specify and calibrate.} All seven predicates have written operational definitions with explicit thresholds, every one fitted on annotator groups 0 to 5 alone and reported on the held-out groups (Chapter 3). The \texttt{near} threshold generalises perfectly to the held-out annotator who used the label (recall 1.00), and the support rules reach held-out F1 0.87.

\textbf{O3, validate.} Per-predicate fidelity against all 8,926 human triplets, against three baselines, with eight ablations and manually audited precision (Chapter 4): mean recall 0.85, and 0.76 on annotators no threshold ever saw. Cluster-bootstrap intervals accompany every headline figure, and an independent validation study re-estimates precision with judges who have no stake in the result (\S4.13).

\textbf{O4, diagnose.} Every one of the 1,689 missed human triplets is attributed to a cause by re-checking rule conditions, so the failure analysis is exhaustive rather than illustrative. Genuine tool error accounts for roughly 7\% of misses; the remainder is calibrated abstention and measured annotator behaviour (\S4.10, \S7.2).

\textbf{O5, test downstream utility.} A controlled experiment trains the same classifier on three label sources under identical conditions (Chapter 5), and the comparison is repeated in the source paper's own benchmark framework with a shared frozen detector and three seeds per arm (Chapter 6).

\textbf{RQ1, accuracy.} Yes, with one qualification. On five of seven predicates the automatic labels match or exceed the human process (0.81 to 1.00 recall, audited precision $\approx$ 1.0), and \texttt{near} is solved outright, resolving the one predicate the source paper reports as failing for every model it benchmarked. The depth pair reaches 0.64/0.66 pooled and 0.84 once two annotator groups' inverted conventions are accounted for; where the tool commits it agrees at 0.95 to 1.00 with six of the seven same-convention annotators, the seventh being the smallest sample in the dataset. The residual limit is monocular ambiguity, which ablation A8 shows a four-times-larger depth model does not fix.

\textbf{RQ2, utility.} Yes, and by a wide margin at this dataset's annotation scale. The auto-trained classifier reaches 0.76 mean recall against held-out human gold versus 0.30 for its human-trained twin, and versus 0.36 when the human labels are stretched by self-training, the standard semi-supervised remedy. In the full benchmark the verdict is conditional and the condition is legible: over three seeds per arm the human-trained model ranks better against human test annotation (mR@100 0.326 against 0.278), the auto-trained model generalises about sixty times better to unseen relation compositions (zR@100 0.172 against 0.003, with disjoint per-seed ranges), and the human advantage is concentrated entirely on the two test annotators carrying measured labelling defects, disappearing against the only one without (0.308 against 0.307).

The single sentence the evidence supports: \textbf{automatic labels are the better supervision wherever ground truth means geometry; human labels remain better wherever ground truth means human annotation practice.} Robot planning needs the former.

\section{Research contributions}

\textbf{For the source dataset and its authors.} An annotator that labels the existing images 20 times more densely in five minutes, in the dataset's own formats, and the written operational definitions the annotation process never had. The fitted \texttt{near} threshold is a direct answer to the future-work request in Wang et al. (2025) for "spatial thresholds for near". Two annotation defects are quantified for the first time: an inverted front/behind convention in two of nine annotator groups, and \texttt{near} used by only three groups with fourfold variation in exhaustiveness.

\textbf{For work on scene-graph benchmarks.} Evidence that ranked recall against sparse, guideline-free annotation partly measures agreement with annotator habits rather than spatial correctness. The evidence is a dissociation: the model that ranks better memorises which pairs annotators record, and the model that generalises to unseen compositions is the one trained on consistent computed labels. The per-annotator decomposition localises the effect precisely, and the seed replication shows which parts of it survive training variance.

\textbf{For weak supervision.} A controlled three-way comparison, on the same features, model, split and seeds, showing that programmatic labels out-teach both scarce human labels and the standard remedy for scarce human labels. The mechanism is measured rather than inferred: self-training contributes roughly a thousand confident negative pseudo-labels for every positive one, propagating the annotators' silence rather than their judgement.

\textbf{Methodologically.} Calibration held out by \textit{annotator} rather than by image; a sparse-gold evaluation protocol that pairs recall with audited precision; exhaustive loss attribution; a way to estimate what annotators would score against one another when they never labelled the same images, by using a deterministic annotator as a fixed common reference (\S4.6); and a reliability check that needs no labels at all, obtained by recovering the fact that an image dataset was cut from a continuous capture and asking whether its labels survive the camera moving (\S4.14). The last of these is the one most likely to transfer. Many robotics datasets are sequences presented as image sets, and wherever that is true, a predicate's agreement with itself across viewpoints separates a rule that is wrong from a rule that is merely uncertain, which is a distinction sparse human annotation cannot draw.

\section{Limitations, and future research and development}

Each limitation below is stated with the specific experiment that would settle it, because that is more useful than an apology.

\textbf{The chain stops short of the robot.} The strongest evidence reaches label quality and model quality; no experiment in this dissertation shows a robot completing more tasks. The instrument exists: 75 prompts across 25 held-out scenes in three conditions (no relations, human relations, automatic relations), with a blind scoring sheet, built and ready to run. This is the single most valuable next experiment, and it is the one that would convert a supported inference into a demonstration.

\textbf{Precision estimates remain partly author-verdicted.} The independent validation study (\S4.13) is deployed and collecting, with coverage stratified so the audit-overlap items reach several raters first. Until it completes, the audited precision figures carry the author's own judgement, conservatively applied and published for checking.

\textbf{One domain, one camera, six object classes.} The fitted constants are dataset-specific by design; what transfers is the calibration procedure. The only out-of-domain evidence is qualitative (\S4.12), where unretuned thresholds behaved correctly on stock video containing objects outside the six classes. A labelled cross-domain sample of a few dozen images would turn that into a measurement, and it is cheap enough that a replication should simply include one.

\textbf{Front/behind is bounded by monocular ambiguity.} Two objects at similar camera distance cannot be ordered from one image, and a larger depth model does not help (A8). The routes out are additional evidence rather than better models: stereo or RGB-D capture, multi-frame structure, or wider surface detection to extend the ground-plane guard, the last of which was built, measured and declined on this dataset (\S4.9.4).

\textbf{Detection bounds full automation.} With a zero-shot detector the end-to-end recall is 0.38, while the relation layer conditional on detection scores 0.85. The gap is detection, not relations, and the source paper's own trained detector (0.93 mAP@50) would close most of it. That check needs only the released weights.

\textbf{Scale is demonstrated without ground truth.} The supervising group supplied the full 2,650-frame capture the released images were cut from, and the pipeline was run over the 1,766 frames nobody has annotated: 562 keyframes after content-adaptive selection, 31 minutes, 185,242 triplets, a predicate distribution 0.032 in total variation from the annotated portion (\S4.15). Capacity and stability on unfamiliar input are therefore measured rather than argued. Correctness on that portion is not, and cannot be without labels; \S4.14's viewpoint-consistency check substitutes self-agreement for truth and should be read as the weaker thing it is. The experiment that would close this is a modest labelled sample from the unannotated frames, a few hundred triplets, which is an afternoon of annotation rather than a research programme.

\textbf{The capture is stereo, and only one eye was used.} The supplied folder is named \texttt{rightimg}, which implies a left counterpart held by the supervising group. True stereo would attack the front/behind bound directly, supplying disparity where this project has only monocular relative depth, and would do so at every frame rather than requiring the robot to move. Multi-frame structure is the weaker fallback if no left channel survives: measured on this sequence, consecutive frames carry 0.08 px of motion, far too little for parallax, with usable baselines appearing only around 20 to 40 frames apart. That route also changes the method's premise, since depth recovered from twenty frames of robot walking is not available to a single-image annotator, so it belongs to a validation instrument rather than to the pipeline.

\section{Personal reflections}

The most useful thing this project taught me was to distrust a number until I know how it was produced.

That lesson arrived early and painfully. My first depth-based results were poor, and the unit tests all passed. The cause was that the dataset's images are stored rotated 180 degrees behind an EXIF flag, so every mask and depth value was being read from an upside-down image while the boxes stayed upright. No automated check could have caught it, because every component was behaving exactly as written. I found it by rendering an image with its boxes drawn on and looking at it. Since then I have rendered and inspected samples after every significant change, and it has caught things twice more.

The second lesson was that "agreement with the humans" is not one target. I began by treating the human labels as ground truth and my disagreements as errors. Measuring them properly showed that three of nine annotator groups used \texttt{near} at all, that two labelled front and behind in opposite directions, and that some labelled support in only one direction. That reframed the project: the interesting question stopped being "how close can I get to the humans" and became "what exactly do the humans mean, and where do they disagree with each other". Almost every later design decision, including calibrating on some annotators and validating on others, follows from that shift.

The third lesson was about my own claims, and it is the one I would most like to have learned sooner. Reading a single benchmark run, I reported that my labels won against the one test annotator with clean conventions. The margin was 0.011 on 73 images. When I later retrained both arms at two more seeds, the arms turned out to be tied, and I withdrew the claim (\S6.3.1). I had been careful about this in one place, fitting thresholds on some annotators and validating on others, and careless about it in another, treating one training run as a result. The discipline was not new to me; I simply had not applied it to model training variance. Recording the retraction in the chapter rather than quietly deleting it felt uncomfortable and is, I think, the right thing to have done.

What I would do differently, given the time again, is front-load the experiments that answer the question the project actually asks. I spent considerable effort proving that the labels are correct and that they train models well, and comparatively little on whether they help a robot decide what to do, which is the entire motivation. The planner experiment should have run in week five rather than sitting built and unrun at the end. I would also have arranged for two people to annotate the same fifty images in week two: that one afternoon would have produced the inter-annotator agreement figure that the whole "comparable to human quality" claim needs, and which I ended up estimating indirectly instead.

Technically, the thing I am most pleased with is not the accuracy figure but the two refinements I built, measured and then declined. Both were plausible, both took real work, and the data said neither helped. Reporting them as null results rather than dropping them was the point at which the project started feeling like research rather than engineering.